% Synthetic fixture: an article the way many authors write one (T1, lmodern, natbib, cleveref).
\documentclass[11pt]{article}
\usepackage[T1]{fontenc}
\usepackage{lmodern}
\usepackage{microtype}
\usepackage{amsmath,amssymb}
\usepackage{graphicx}
\usepackage{caption}
\usepackage{subcaption}
\usepackage{siunitx}
\usepackage[numbers,sort&compress]{natbib}
\usepackage[colorlinks=true,allcolors=blue]{hyperref}
\usepackage[capitalise]{cleveref}

\newcommand{\Rate}{\lambda}

\title{Tide Gauges Along a Winding Coast}
\author{R. Okafor \and M. Lindqvist}
\date{}

\begin{document}
\maketitle

\begin{abstract}
We compare three ways of reading a tide gauge that sits in a narrow inlet, and
show that a simple smoothing rule matches the harbour master's tables.
\end{abstract}

\section{Introduction}\label{sec:intro}
Our gauge records the water level every \SI{6}{\minute}, which gives
\num{240} readings a day. The earlier survey \citep[see][p.~3]{harbour94} used
hourly values; later work \citealp{inlet07,quay12} kept that habit.

The raw logs are kept at \url{https://example.org/tides/raw%20logs/2024#march}
and the notes at \href{https://example.org/tides/notes%20v2}{the project page}.\footnote{The
page is mirrored; see the archive in \cref{sec:method}.}

{\bf Stations.} The three stations are listed below.
\begin{description}
  \item[North] A pier gauge, read by hand until 1998.
  \item[Middle] A float gauge in the stilling well.
  \item[South] A pressure sensor, new this year.
\end{description}

\section{Method}\label{sec:method}
We take the level $h(t)$ and write
\begin{equation}\label{eq:level}
  h(t) = a_0 + \sum_{k=1}^{4} a_k \cos(\omega_k t + \phi_k).
\end{equation}

Three steps follow:
\begin{enumerate}
  \item Remove spikes, in two parts:
    \begin{itemize}
      \item values beyond \SI{3}{\metre};
      \item steps larger than \SI{0.5}{\metre} in one reading.
    \end{itemize}
  \item Fit \cref{eq:level} by least squares.
  \item Smooth the residual with a window of \num{5} readings.
\end{enumerate}

The fit follows \citep{harbour94}; an alternative is described by
\citealp{inlet07}. Pressure sensors drift, so the south gauge is checked
against the tables each month.


\section{Results}\label{sec:results}
\begin{figure}[ht]
  \centering
  \begin{subfigure}[b]{0.45\textwidth}
    \centering
    \fbox{\parbox{0.8\textwidth}{\centering raw trace}}
    \caption{Raw readings.}\label{fig:raw}
  \end{subfigure}\hfill
  \begin{subfigure}[b]{0.45\textwidth}
    \centering
    \fbox{\parbox{0.8\textwidth}{\centering smoothed trace}}
    \caption{After smoothing.}\label{fig:smooth}
  \end{subfigure}
  \caption{Two views of one week of readings.}\label{fig:both}
\end{figure}

\Cref{fig:raw,fig:smooth} show the traces side by side, and \cref{fig:both} gives
the pair. See also \crefrange{eq:level}{eq:rate} and \citet{quay12}.

The level is zero when $x \text{ if $y$ holds}$ and the rate is positive otherwise.
The mean over the \SI{12.4}{\hour} cycle is
\begin{equation}\label{eq:rate}
  \Rate = \frac{1}{T}\int_0^T \dot{h}(t)\,dt .
\end{equation}

\section{Conclusion}
A short smoothing window is enough. The data are open, and the code is small.

\bibliographystyle{plainnat}
\bibliography{refs}
\end{document}
